\chapter{Extended Background} \label{cha:extendBackground}
\par\vspace{1em}
\section{Literature Review}

\subsection{Fundamentals of LSB Steganography}
LSB embedding stands out as one of the most popular and accessible steganographic techniques. 
LSB steganography replaces the value of a least significant bit of a pixel with the bit from the hidden message, as described by \textcite{Chan2004}. 
For instance, in an 8-bit grayscale image, modifying the least significant bit of a pixel value from 10110110 (182) to 10110111 (183) results in only one change in intensity level (see figure 2.1), which is essentially undetectable by the human eye. 
The advantage of LSB insertion lies in its simplicity and high embedding capacity; however, these advantages also exhibit statistical abnormalities that can be utilized in the detection process.
The image format used to contain the hidden information plays a vital role in both the effectiveness of the steganographic embedding and its detectability. 
While not the only lossless format, PNG (Portable Network Graphics) images are highly relevant to LSB steganography given their characteristics as a lossless compressible file format. 
While JPEG is a useful format, it is lossy and often destroys hidden data when compression artifacts are introduced \cite{Elgabar2013}. 
PNG images and their deflate compression algorithm preserve pixel values exactly, thus maintaining an ideal medium for LSB steganography because the image maintains the exact bit values required for message recovery.

\begin{figure}[ht]
    \centering
    \includegraphics[width=0.8\textwidth]{lsb_illustration.png}
    \caption{LSB Illustration.}
    \label{fig:roc_matching}
\end{figure}


\subsection{Traditional Steganalysis}
\textcite{Fridrich2001} introduced the RS (Regular-Singular) steganalysis method to detect 
LSB steganography in grayscale images. Unlike histogram attacks (like Chi-Square) 
that look for global frequency anomalies, RS analysis uses spatial correlations between 
groups of pixels to estimate the initial cover image.

RS analysis works by dividing the image into disjoint groups of $n$ 
adjacent pixels (e.g., $2 \times 2$ blocks). A discrimination function, $f$, 
is defined to measure the "noisiness" or variation within each group. As defined by 
\textcite{Fridrich2001}, for a group of pixels $G = (x_1, x_2, \dots, x_n)$, 
the discrimination function is calculated as the sum of absolute differences between adjacent pixels:

\begin{equation}
    f(x_1, x_2, \dots, x_n) = \sum_{i=1}^{n-1} |x_{i+1} - x_i|
\end{equation}

By applying a flipping operation (mask $M$) to the LSBs of the pixels in a group, the block's noisiness changes. Based on how the discrimination function value changes after flipping (denoted as $F_M(G)$), groups are categorized into:

\begin{itemize}
    \item \textbf{Regular (R):} Flipping increases the noisiness ($f(F_M(G)) > f(G)$).
    \item \textbf{Singular (S):} Flipping decreases the noisiness ($f(F_M(G)) < f(G)$).
    \item \textbf{Unusable (U):} Flipping does not change the noisiness ($f(F_M(G)) = f(G)$).
\end{itemize}

Statistical symmetry is the core premise of RS analysis. 
In an image with no embedding, the number of Regular Groups ($R_M$) should be approximately equal to the 
number of Regular groups obtained using an inverse mask ($R_{-M}$); the case is similar for Singular groups:

\begin{equation}
    R_M \approx R_{-M} \quad \text{and} \quad S_M \approx S_{-M}
\end{equation}


The natural balance is disturbed by LSB embedding, and the disturbance grows as the payload increases. 
As the size of this modified payload increases, the difference between ($R_M$) and ($S_M$) will become smaller. 
This process continues until the image becomes completely randomized; at that point, $R_M$ and $S_M$ converge. 
Since RS Analysis measures absolute deviation from this perfect symmetry, 
it not only allows for detection of steganography but also allows for a 
more accurate estimation of embedded message length than earlier first-order statistical tests.

\subsection{Deep Learning in Steganalysis}
The limitations of handcrafted statistical features, such as the rigid assumptions regarding spatial correlation in RS-analysis, paved the way for deep learning approaches. 
While RS-analysis relies on explicit knowledge of the embedding process, deep learning (DL) methods can automatically learn complex patterns that distinguish cover images from stego-images. 
This capability gives DL methods an advantage, as a Convolutional Neural Network (CNN) can learn to identify steganographic artifacts from raw data that remain statistically invisible to fixed-model testing methods \cite{Farooq2024}.

The field has experienced a paradigm shift from spatial-domain analysis toward adaptive domain techniques driven by artificial intelligence \cite{Apau2024}. 
The primary method that has emerged is the use of CNNs. 
However, the implementation of CNNs for steganalysis provides unique challenges compared to traditional computer vision tasks. 
In a classic example of image recognition (e.g., recognizing a cat in a photo), the object being identified is the dominant signal.

Steganographic modifications, conversely, are designed to be imperceptible, hidden behind the
"strong" noise of the image content \cite{Qian2015}. 
Standard CNN techniques usually fail to detect these weak signals because the high-amplitude content suppresses them. 
To address this, \textcite{Qian2015} proposed a customized deep learning model known as the Gaussian-Neuron CNN (GNCNN).

The GNCNN introduces a pre-processing layer designed to suppress image content and enhance the signal-to-noise ratio of the steganographic embedding. 
A high-pass filter is applied before the data enters the convolutional layers, effectively removing the "cover image" and leaving the "noise residual" where steganographic artifacts reside. 
This is similar to the residual calculation step in traditional steganalysis but is integrated directly into the learning pipeline.

The key benefit of using CNNs over traditional statistical methods, such as RS-analysis, is that the CNN architecture combines both feature extraction and classification into one unified structure. 
In traditional methods, feature extraction is completed first, and those features are used to train the classifier independently; with CNNs, both feature extraction and classification occur in concert using backpropagation \cite{Qian2015}.

\subsection{Sensitivity and Robustness}
Researching payload size raises issues regarding detection sensitivity. 
RS-Analysis relies on comparing variations within block statistics (based on $R_M$ versus $S_M$). 
As the payload decreases, the difference becomes negligible. At low payload size (below 0.05 bpp), the modification to pixel groups does not create enough signal to differentiate the image from natural noise. 
While statistical approaches struggle to create an accurate threshold in these low-payload situations, deep learning models have shown superior performance \cite{Apau2024}.

Combating the false positive rate is also vital for operational security. 
A key weakness of statistical tests is the occurrence of false alarms when viewing noisy images or those with distinct textures. 
Scanned photographs or images altered using specific filters often possess statistical properties that confuse RS steganalysis, incorrectly classifying clean images as stego-images \cite{Fridrich2002}. 
CNNs typically offer better generalizability. After being trained on large datasets, deep learning models can learn to detect the specific artificial noise introduced by LSB embedding, reducing the false positive rate compared to rigid statistical tests \cite{Qian2015} \cite{Apau2024}. 
